\section{O`zgarmas tokning ishi va quvvati. Joul-Lens qonuni}


\begin{enumerate}

\item Berilgan formulalardan qaysi biri Joul–Lens qonunini ifodalaydi? 
\begin{enumerate}[label=\Alph*)]
\item $Q = I^2Rt$
\item $Q = \lambda m$
\item $Q = cm(t_2 - t_1)$
\item $Q = A + \Delta U$
\end{enumerate}

\item Qarshiligi $2~\Omega$ bo‘lgan simdan $15$ A tok o‘tadi. $2$ minut ichida unda qancha energiya ajraladi (kJ)? 
\begin{enumerate}[label=\Alph*)]
\item 27
\item 54
\item 450
\item 900
\end{enumerate}

\item $50~\Omega$ qarshilikli o‘tkazgich orqali $10$ minut davomida qanday tok o‘tkazilganda, $120$ kJ issiqlik ajralib chiqadi (A)?
\begin{enumerate}[label=\Alph*)]
\item 2
\item 3
\item 4
\item 5
\end{enumerate}

\item Qarshiligi $12~\Omega$ bo‘lgan o‘tkazgich $120$ V kuchlanishli tarmoqqa ulandi. O‘tkazgichda $10$ min ichida qancha issiqlik ajraladi?
\begin{enumerate}[label=\Alph*)]
\item 72 kJ
\item 0,72 kJ
\item 720 kJ
\item 7,2 MJ
\end{enumerate}

\item Agar $R_1 = R_2 = R_3$ bo‘lsa, rasmda ko‘rsatilgan rezistorlardan bir xil vaqt davomida ajralib chiqayotgan issiqlik miqdorlarining nisbati $Q_1 : Q_2 : Q_3$ ni aniqlang.

\begin{circuitikz}[european]
\draw
  % Kirish simi va R1 rezistori (ketma-ket qism)
  (0,0) to[R=$R_1$, *-] (2,0)
  
  % Parallel tarmoqlanish boshlanishi
  (2,0) -- (2,0.7) to[R=$R_2$] (4,0.7) -- (4,0)
  (2,0) -- (2,-0.7) to[R=$R_3$] (4,-0.7) -- (4,0)
  
  % Chiqish simi va tok kuchi I
  (4,0) -- (5.5,0)
  (4.5,0) node[currarrow, label=above:$I$] {};
\end{circuitikz}

\begin{enumerate}[label=\Alph*)]
\item 4:1:1
\item 1:1:1
\item 4:4:1
\item 1:4:4
\end{enumerate}

\item $R$ va $2R$ qarshiliklar o‘zgarmas kuchlanish manbaiga parallel ulangan. $R$ qarshilikda $100$ J issiqlik miqdori ajralganda, $2R$ qarshilikda qancha issiqlik miqdori ajraladi (J)?
\begin{enumerate}[label=\Alph*)]
\item 400
\item 200
\item 100
\item 50
\end{enumerate}


\item Elektr choynakning ikki chulg‘ami bor. Ulardan biri manbaga ulanganda choynakdagi suv $30$ minutda, ikkinchisi ulanganda esa $20$ minutda qaynaydi. Agar ikkala chulg‘am manbaga ketma-ket ulansa, choynakdagi suv necha minutda qaynaydi?
\begin{enumerate}[label=\Alph*)]
\item 25
\item 30
\item 40
\item 50
\end{enumerate}

\item Elektr isitkich bir xil $R$ qarshilikli uchta spiralga ega. Agar uchala spiral parallel ulansa, unga qo‘yilgan idishdagi suv $6$ minutda qaynaydi. Spirallar rasmda ko‘rsatilgandek ulansa (ikki spiral parallel va biri ularga ketma-ket), shu suv necha minutda qaynaydi?
\begin{enumerate}[label=\Alph*)]
\item 24
\item 4
\item 12
\item 16
\end{enumerate}


\item Chulg‘amning qarshiligi $50~\Omega$ bo‘lgan elektr choynakda temperaturasi $0^\circ$C bo‘lgan $600$~cm$^3$ suv bor. Agar tarmoqning kuchlanishi $200$~V, choynakning FIK $60\%$ bo‘lsa, undagi hamma suvni qaynatib, bug‘ga aylantirish uchun necha minut vaqt kerak bo‘ladi? Suv uchun $c=4,2$~kJ/(kg$\cdot$K), $r=2,3$~MJ/kg.
\begin{enumerate}[label=\Alph*)]
\item 57
\item 28
\item 32
\item 41
\end{enumerate}

\item Qaysi qarshilikda eng ko‘p issiqlik ajraladi? $R_1=1~\Omega, R_2=R_3=2~\Omega, R_4=4~\Omega$. (Rasmga qarang)

\begin{circuitikz}[european]
\draw
  % Kirish simi va tugun nuqtasi
  (0,0) -- (1,0) to[short, *-] (1,0)
  
  % Yuqori tarmoq: R1 va R2 ketma-ket
  (1,0) -- (1,1) to[R=$R_1$] (3,1) to[R=$R_2$] (5,1) -- (5,0)
  
  % Pastki tarmoq: R3 va R4 ketma-ket
  (1,0) -- (1,-1) to[R=$R_3$] (3,-1) to[R=$R_4$] (5,-1) -- (5,0)
  
  % Chiqish simi va tugun nuqtasi
  (5,0) to[short, -*] (6,0);
\end{circuitikz}


\begin{enumerate}[label=\Alph*)]
\item $R_1$ da
\item $R_2$ da
\item $R_3$ da
\item $R_4$ da
\end{enumerate}

\item Qarshiligi $R_1$ bo‘lgan qaynatgich biror miqdor suvni $t_1$ vaqtda qaynatadi. Qarshiligi $R_2$ bo‘lgan qaynatgich esa shu sharoitda o‘sha suvni $t_2$ vaqtda qaynatadi. Ikkala qaynatgich parallel ulansa o‘sha suv qanday vaqtda qaynaydi?
\begin{enumerate}[label=\Alph*)]
\item $t_1+t_2$
\item $\frac{t_1 t_2}{t_1+t_2}$
\item $\sqrt{t_1 t_2}$
\item $\frac{t_1+t_2}{2}$
\end{enumerate}

\item Elektr choynakda ikkita chulg‘am bor. Birinchi chulg‘am ulanganda, suv $t_1$ vaqtda, ikkinchisi ulanganda, $t_2$ vaqtda qaynaydi. Ikkala chulg‘am ketma-ket ulanganda qaynash vaqti, ular parallel ulangandagi qaynash vaqtidan necha marta katta bo‘ladi?
\begin{enumerate}[label=\Alph*)]
\item $\frac{2t_1 t_2}{t_1+t_2}$
\item $\frac{t_1+t_2}{t_1 t_2}$
\item $\frac{(t_1+t_2)^2}{t_1 t_2}$
\item $\frac{t_1-t_2}{t_1+t_2}$
\end{enumerate}

\item Ikkita bir xil elektr plita parallel ulansa, suv $t$ vaqtda qaynaydi. Shu plitalar ketma-ket ulansa, shu suv qancha vaqtda qaynaydi?
\begin{enumerate}[label=\Alph*)]
\item $2t$
\item $t/2$
\item $t/3$
\item $4t$
\end{enumerate}


\item Elektr choynakning ikkita chulg‘ami bor. Ulardan biri manbaga ulanganda choynakdagi suv $30$ minutda, ikkinchisi ulanganda esa $20$ minutda qaynaydi. Agar ikkala chulg‘am manbaga ketma-ket ulansa, choynakdagi suv necha minutda qaynaydi?
\begin{enumerate}[label=\Alph*)]
\item 25
\item 30
\item 40
\item 50
\end{enumerate}

\item Elektr isitkich bir xil $R$ qarshilikli uchta tarmoqqa ega. Agar har uchala tarmoq o‘zaro parallel ulansa, unga qo‘yilgan idishdagi suv $6$ min da qaynaydi. Tarmoqlar o‘zaro ketma-ket ulanganida shu suv necha minutda qaynar edi?
\begin{enumerate}[label=\Alph*)]
\item 60
\item 54
\item 44
\item 38
\end{enumerate}

\item Elektr isitkich bir xil $R$ qarshilikli uchta spiralga ega. Agar uchala spiral parallel ulansa, unga qo‘yilgan idishdagi suv $6$ minutda qaynaydi. Spirallar rasmda ko‘rsatilgandek ulansa, shu suv necha minutda qaynaydi?

\begin{circuitikz}[european]
\draw
  % Kirish tugun nuqtasi
  (0,0) to[short, *-] (0.5,0)
  
  % Yuqori tarmoq: ikkita R rezistor ketma-ket
  (0.5,0) -- (0.5,0.8) to[R=R] (2.5,0.8) to[R=R] (4.5,0.8) -- (4.5,0)
  
  % Pastki tarmoq: bitta R rezistor
  (0.5,0) -- (0.5,-0.8) to[R=R] (4.5,-0.8) -- (4.5,0)
  
  % Chiqish tugun nuqtasi
  (4.5,0) to[short, -*] (5,0);
\end{circuitikz}

\begin{enumerate}[label=\Alph*)]
\item 24
\item 4
\item 12
\item 16
\end{enumerate}

\item Elektr plita suvni $10$ minutda qaynatadi. Spiralning bir qismi uzilgandan so‘ng, shu suvni $12$ min da qaynatdi. Spiral dastlabki uzunligining qanday qismiga kaltalashgan? Tok kuchini o‘zgarmas deb hisoblang.
\begin{enumerate}[label=\Alph*)]
\item $1/3$
\item $1/8$
\item $1/12$
\item $1/6$
\end{enumerate}

\item Suv ikkita bir xil spiral bilan isitiladi. Spirallar dastlab o‘zgarmas kuchlanish manbayiga $t$ vaqt davomida parallel, keyin xuddi shuncha vaqt ketma-ket ulandi. Agar dastlabki $t$ vaqtda suvning temperaturasi $10^\circ$C ga ortgan bo‘lsa, $2t$ umumiy vaqtda qanchaga ($^\circ$C) ortgan?
\begin{enumerate}[label=\Alph*)]
\item 20
\item 15
\item 17,5
\item 12,5
\end{enumerate}




\end{enumerate}



\begin{table}[!ht]
\centering
\begin{tabular}{|c|c|c|c|c|c|c|c|c|c|}
\hline
1a & 2d & 3b & 4a & 5d & 6c & 7d & 8c & 9a & 10a \\ \hline
11a & 12d & 13b & 14a & 15b & 16b & 17c & 18d & & \\ \hline
\end{tabular}
\end{table}


